\section{Future Plans}\label{sec:future}

\subsection{Promotion and adoption plan}\label{sec:promo}

The path to reach the target population runs through the institutions the
cluster already trusts rather than through direct marketing to individual
factories. TIFI's $\sim$700 member firms and MIRDC's subsidized counseling
channel are the primary distribution routes; a single verified Gangshan pilot,
documented, becomes the reference case that the association can propagate. In
parallel, the project seeks integration of its computation layer into MOENV's
existing CBAM help platform and SMESA's CAAS service network\cite{tpt-cbam,caas},
so that CarbonPass reaches SMEs as an extension of services they already
consult, at the moment of need. Public visibility from the competition itself,
and the openly published datasets (Section~\ref{sec:opendata}), further lower
the cost of discovery.

\subsection{Engaged collaboration units}\label{sec:engaged}

At submission, collaboration is at the outreach stage and is stated honestly as
such: TIFI member firms, MIRDC, and Kaohsiung EPB's voluntary-reduction
advisory team are being approached as pilot and validation partners for the
mentorship period. The competition's own public--private matchmaking and
field-validation support are expected to formalize these relationships for
advancing teams\cite{handbook}.

\subsection{Future planned collaboration units}\label{sec:planned}

\begin{table}[htbp]
\caption{Planned collaboration and the expansion sequence}
\label{tab:expansion}
\small
\begin{tabularx}{\textwidth}{@{}p{0.24\textwidth}p{0.28\textwidth}X@{}}
\toprule
\textbf{Horizon} & \textbf{Partners} & \textbf{What happens} \\
\midrule
2026--2027 (Taiwan beachhead) & TIFI, MIRDC, MOENV, SMESA/MOEA & Harden the pilot into cluster-scale rollout; propose integration into the MOENV CBAM platform and CAAS network; publish open datasets and open-source components\cite{tpt-cbam,caas} \\
2027--2028 (regulatory expansion) & EU importers; standards bodies & If the European Parliament confirms the downstream extension, onboard newly covered categories --- springs, wire, hand tools, household metal articles --- from the same engine\cite{akin}; extend outputs toward Digital Product Passport payloads\cite{espr} \\
2027 onward (domestic breadth) & MOENV; FSC (green finance); banks & Serve SMEs entering Taiwan's lowered carbon-fee threshold and ETS pilot era\cite{ecobiz}; provide bank-ready carbon data for sustainability-linked lending \\
2028+ (international) & SME-exporting economies in Asia & Port the architecture to Vietnam, Thailand, Malaysia, and India, which face the same mechanism with thinner tooling \\
\bottomrule
\end{tabularx}
\end{table}

The expansion is one engine with a widening catalogue: fasteners (2026, live
compulsion) $\to$ downstream steel and aluminium categories (2028, pending)
$\to$ Digital Product Passport sectors (2027--2030) $\to$ domestic carbon-fee
entrants (2027 onward). The same open-data grid signal, local-first document
AI, and EU-format output serve every stage.

\begin{spotlight}{Beyond Taiwan.}
The carbon-data wall is not Taiwan's alone. Every SME-exporting economy faces
the same EU mechanism, most with weaker digital infrastructure than Taiwan's.
CarbonPass is portable by design --- swap the national grid feed, emission
coefficients, and tariff tables, and the same architecture serves another
country's exporters. A Taiwanese answer to a global problem: digital inclusion
delivered as trade capability.
\end{spotlight}
